\addtocontents{toc}{~\hfill\textbf{Page}\par}
\addcontentsline{toc}{section}{Abstract}
\textbf{Abstract}\\

Large Vision-Language Models (LVLMs) have achieved strong performance across medical imaging tasks, yet they remain prone to factual inconsistencies, poor visual grounding, and misalignment with clinically meaningful feedback. Among the post-training remedies proposed for these failures, Direct Preference Optimization (DPO) has been widely adopted for its ability to learn from comparisons without an explicit reward model. This thesis argues that its standard formulation is nonetheless structurally mismatched to clinical use in three ways. First, its sequence-level reward treats a diagnostically decisive span such as an anatomical side or a lesion attribute exactly like generic filler text, so the training signal is dominated by tokens that carry no clinical meaning. Second, because expert preference annotation is expensive, practitioners substitute the supervised fine-tuning reference as the preferred response; the resulting off-policy distribution shift lets the model raise its preference margin by imitating the reference's brevity and formatting rather than by correcting the clinical error, a failure mode we characterize as \emph{stylistic reward hacking}. Third, existing objectives contain no explicit visual grounding constraint, leaving models insensitive to the subtle pathological features that determine a diagnosis.

This thesis introduces \ourloss{}, a fine-grained, on-policy alignment framework built around three corresponding components. Preference pairs are constructed by minimally editing the model's own generations, correcting only clinically erroneous spans while preserving the original linguistic style, so the learning signal isolates factuality from surface form. A bidirectional token-wise KL regularizer balances the mean-seeking behaviour of forward KL, which preserves the base model's linguistic coverage, against the mode-seeking tail suppression of reverse KL, which is what prevents spuriously fluent but incorrect tokens. A visual-contrastive grounding objective pairs each image with a lesion-corrupted counterpart and reverses the preference direction when the supporting evidence has been removed, penalizing answers that are correct for the wrong reasons.

Across two state-of-the-art backbones and three clinical benchmarks spanning visual question answering and radiology report generation, \ourloss{} improves on both standard DPO and prior fine-grained alignment methods, with an average relative gain of $10.24\%$. Layer-wise attention analysis shows that the aligned models attend more consistently to diagnostically relevant image regions, and controlled ablations isolate the contribution of each component: the bidirectional regularizer, lesion-targeted rather than global image corruption, and on-policy rather than reference-style preference pairs.
